\documentclass[12pt,a4paper]{article}
\usepackage[T1]{fontenc}
\usepackage[utf8]{inputenc}
\usepackage{polski}
\usepackage{amsmath}
\usepackage{geometry}
\geometry{a4paper, margin=2.5cm}
\usepackage{enumitem}
\usepackage{titlesec}

\usepackage{hyperref}
% Automatyczna numeracja sekcji z paragrafem (§)
\renewcommand{\thesection}{\arabic{section}}

% Wyśrodkowanie tytułów sekcji
\titleformat{\section}
  {\normalfont\large\bfseries\centering}{\S~\thesection.}{1em}{}

\title{\huge Regulamin \textbf{Smartfonowych Regat}\\
\Large o Puchar Honorowego Komandora AKŻ AGH\\
kpt. Leszka Kurcza}

\date{}
\author{}

% Zmiana numeracji punktów na format paragraf.punkt.podpunkt
\setlist[enumerate,1]{label=\thesection.\arabic{enumi}.}
\setlist[enumerate,2]{label=\thesection.\arabic{enumi}.\arabic{enumii}.}
\setlist[enumerate,3]{label=\alph{enumiii})}

\begin{document}

\maketitle


% \section{Wstęp}
% \begin{enumerate}
%     \item a
% \end{enumerate}

\section{Postanowienia ogólne}
\begin{enumerate}
    \item Regulamin dotyczy Smartfonowych Regat o Puchar Honorowego Komandora AKŻ AGH, kpt. Leszka Kurcza, zwanych dalej \textit{,,Regatami''}.
    \item Organizatorem Regat jest \textbf{Akademicki Klub Żeglarski AGH}.
    \item Regaty odbywają się co roku podczas \textbf{Zakończenia Sezonu AKŻ AGH} lub w innym terminie ustalonym przez Organizatora.
    \item Uczestnikami Regat mogą być uczestnicy Zakończenia Sezonu oraz osoby dopuszczone przez Organizatora.
    \item Nad przebiegiem Regat czuwa \textbf{Sędzia}, niebędący Uczestnikiem.
    W razie jego nieobecności -- Komisja Sędziowska, składająca się z minimum trzech uczestników Regat. Oba określane dalej jako \textit{,,Sędzia''}.
    \item Regaty mają charakter \textbf{towarzyski}, uczestnicy zobowiązani są do polubownego rozstrzygania sporów z poszanowaniem etykiety żeglarskiej.
    \item Sędzia rozstrzyga wszelkie kwestie, które nie zostały uwzględnione w Regulaminie.
\end{enumerate}


\section{Zasady żeglugi}
\begin{enumerate}
    \item Podczas całych regat obowiązuje prawo drogi na wodach śródlądowych. Wszelkie odstępstwa od tego są przez Sędziego przekazywane na Odprawie.
    \item Dodatkowo, w miejscach ustalonych przez Regulamin i Sędziego stosuje się Regatowe Prawo Drogi -- skrót dostępny w \ref{sec:RegatowePrawoDrogi}.
    \item Regaty odbywają się na Jeziorze Rożnowskim na jachtach klasy Omega.
    \item Załoga powinna liczyć \textbf{cztery osoby}, z których minimum jedna posiada odpowiednie uprawnienia żeglarskie (zwana dalej \textit{Sternikiem}).
    \item Szczegółowe zasady dotyczące przebiegu i harmonogramu regat w danym dniu podejmowane są przez Sędziego i przekazywane w trakcie Odprawy przed rozpoczęciem biegów.
    \item Wszelkie informacje o zmianach w przebiegu regat, w szczególności spowodowanych względami bezpieczeństwa, będą przekazywane na bieżąco przez Sędziego.
    \item Sędzia decyduje o konieczności założenia przez załogi środków ratunkowych. W przypadku niedostosowania się do decyzji Sędziego, załoga może zostać zdyskwalifikowana. 
\end{enumerate}


\section{Przebieg regat}
\begin{enumerate}
    \item Regaty są rozgrywane we \textit{Flotach} (grupach załóg):
    \begin{enumerate}
        \item Maksymalna liczba załóg we \textit{Flocie} ustalana jest przez Sędziego, zależnie od liczby dostępnych jachtów i warunków na akwenie.
        \item Punkty są przydzielane wewnątrz każdej \textit{Floty} niezależnie.
        \item W danej \textit{Części} regat, \textit{Floty} kolejno rozgrywają biegi.
        \item \textit{Seria} biegów składa się z jednego biegu każdej \textit{Floty}.
    \end{enumerate}
    \item Regaty składają się z dwóch \textit{Części}:
    \begin{enumerate}
        \item Eliminacje:
        \begin{enumerate}
            \item Formuła \textbf{Klasyczna}: załogi pokonują trasę wyznaczoną przez Sędziego, od startu do mety, o klasyfikacji decyduje czas pokonania trasy.
            \item Czas jest mierzony przez Sędziego lub wyznaczoną Komisję Pomiarową.
            \item W eliminacjach załogi są podzielone na \textit{Floty} na podstawie losowania lub innego kryterium ustalonego przez Sędziego.
            \item Obowiązuje prawo drogi na wodach śródlądowych oraz Regatowe Prawo Drogi przedstawione Wszelkie odstępstwa od tego są przez Sędziego przekazywane na Odprawie. 
        \end{enumerate}
        \item Finały:
        \begin{enumerate}
            \item Formuła \textbf{Smartfonowa} (tradycja AKŻ AGH): załogi pokonują dowolną trasę, w ograniczonym czasie określonym przez Sędziego, o klasyfikacji decyduje długość pokonanej trasy.
            \item Przebieg trasy jest mierzony przez załogi za pomocą Aplikacji.
            \item W finałach załogi są podzielone na \textit{Floty} na podstawie wyników klasyfikacji z eliminacji (opisane w \ref{punktacja:eliminacje}).
            \item Floty w finale dzielą się na złotą, srebrną itd. (przyjęte nazwy: złota-,,Wilki   Morskie'', kolejne np. ,,Węże Rzeczne'', ,,Szczury Lądowe''). Liczba flot na dany rok ustalana jest na podstawie liczby Uczestników.
            \item Obowiązuje prawo drogi na wodach śródlądowych. Wszelkie odstępstwa od tego są przez Sędziego przekazywane na Odprawie.
        \end{enumerate}
    \end{enumerate}
    
    \item Komunikacja i wszystkie sprawy organizacyjne odbywają się podczas \textit{Odpraw} Sterników:
    \begin{enumerate}
        \item Odprawy prowadzi Sędzia.
        \item Na Odprawach obowiązkowa jest obecność każdego Sternika (lub przedstawiciela załogi). Brak obecności skutkuje klasyfikacją załogi na ostatnim miejscu w następnym biegu.
        \item Odprawy odbywają się przed, pomiędzy oraz po każdej \textit{Serii} biegów.
        \item Na pierwszej Odprawie Sędzia przekazuje szczegółowe informacje dotyczące przebiegu regat, w tym zasady bezpieczeństwa, zasady żeglugi, podział na \textit{Floty}, czas trwania biegów, zasady punktacji oraz inne istotne informacje.
    \end{enumerate}

    \item Przed każdym biegiem załogi mają obowiązek stawienia się na keji (Gotowość), w gotowości do przejęcia jachtu. Wtedy też Sędzia może zadecydować o opóźnieniu startu następnego biegu i ustalić późniejszą godzinę Gotowości.
    \item W szczególnym przypadku (złe warunki pogodowe) Sędzia wraz ze wszystkimi Sternikami może ustalić zastępczą metodę klasyfikacji dla danego biegu.
    \item Podczas każdego biegu załogi wykonują pomiary za pomocą \textbf{aplikacji mobilnej} -- Strava, Żegluj lub inna, ustalona i zaakceptowana przez Sędziego. Zwana dalej \textit{,,Aplikacją''}.
    \item Aplikacja powinna działać na ustawieniach domyślnych, a użytkownicy powinni posiadać zarejestrowane konta w aplikacji. Powinna też być umożliwiona funkcja publikowania na platformach społecznościowych (np. Facebook) wyników po zakończeniu wyścigu.
    \item Po zakończeniu każdego biegu, w ciągu 15 minut (lub inny czas ustalony przez Sędziego), wynik z Aplikacji powinien znaleźć się na FB, opisany odpowiednimi tzw. “hasztagami”, np. \texttt{\#akzagh}, \texttt{\#AKZ}, lub innymi podanymi przez organizatora,
    \item Korzystamy ze sprzętu dostępnego nad akwenem, dlatego też regaty zawsze mają charakter \textbf{przesiadkowy}.
\end{enumerate}



\section{Punktacja, klasyfikacja}
\begin{enumerate}
    \item Miejsca we wszystkich biegach punktowane są zgodnie z numerem miejsca (tj. 1. miejsce -- 1pkt, 2. miejsce -- 2pkt itd.).
    \item Liczba biegów w każdej \textit{Części} regat ustalana jest przez Sędziego przed Regatami.
    \item W obu \textit{Częściach} regat wygrywa załoga z \textbf{najmniejszą liczbą punktów}, pozostałe załogi klasyfikowane są w odpowiedniej kolejności. 
    \item Klasyfikacja eliminacyjna (formuła Klasyczna): \label{punktacja:eliminacje}
    \begin{enumerate}
        \item W danym biegu, w obrębie danej \textit{Floty}, załogi zajmują miejsca w zależności od \textbf{czasu} przybycia do mety.
        \item Klasyfikacja w eliminacjach decyduje o przydzieleniu załóg do \textit{Flot} w finale.
        \item W przypadku remisu punktowego, decyduje wartość procentowego odchylenia czasu od czasu zwycięzcy w danym biegu (mniejsza wartość jest lepsza):
        \begin{equation*}
        \text{odchylenie \%} = \frac{\text{CzasZałogi} - \text{CzasZwycięzcy}}{\text{CzasZwycięzcy}} \times 100\%
        \end{equation*}
        \item W przypadku kilku \textit{Serii} biegów, decyduje suma punktów (procentów odchylenia).
    \end{enumerate}
    \item Klasyfikacja Finałowa (formuła Smartfonowa):
    \begin{enumerate}
        \item Klasyfikacja finałowa jest na podstawie kategorii \textbf{Głównej (Sztandarowej)}: Najdłuższa przebyta \textbf{droga} w określonym czasie.
        \item Klasyfikacji załóg odbywa się na podstawie wyników zebranych w Aplikacji.
        \item Każda załoga zobowiązana jest do posiadania \textbf{minimum dwóch urządzeń} z zainstalowaną tą aplikacją.
        \item Przy zbieraniu wyników z dwóch aplikacji liczy się gorszy wynik.
        \item W każdej \textit{Flocie} załogi zajmują miejsca w zależności od długości przebytej trasy.
    \end{enumerate}
    \item Kategorie dodatkowe:
    \begin{enumerate}
        \item ,,Najszybsza załoga'' -- Najszybsza zarejestrowana prędkość
        \item ,,Mistrzowie Nawigacji'' -- Najbardziej kreatywna trasa. Przyznawana przez Honorowego Komandora lub, w przypadku jego braku, przez Sędziego lub Komisję po zakończeniu Regat.
    \end{enumerate}
    \item Czas trwania każdego biegu, godzina startu oraz zakończenia jest podana przez Sędziego w Harmonogramie. W tym czasie załoganci, pływając po dowolnym rejonie Jeziora (ograniczonym przez Sędziego), żeglują swobodnie po jeziorze, mając uruchomioną aplikację. Przekroczenie czasu skutkuje karami.
\end{enumerate}



\section{Nagrody}
\begin{enumerate}
    \item Zwycięzcy kategorii Głównej, we flocie złotej (,,Wilki Morskie''), zdobywają \textbf{Puchar Przechodni} oraz mają zagwarantowane umieszczenie nazwisk załogi na tym trofeum.
    \item Zwycięzcy kategorii Głównej w pozostałych flotach finałowych zdobywają przewidziane na dany rok przez Organizatora nagrody.
    \item Zwycięzcy pozostałych kategorii otrzymują przewidziane na dany rok przez Organizatora nagrody.
    \item Wszystkim zwycięzcom przysługuje chwała i miejsce w historii Akademickiego Klubu Żeglarskiego AGH. 
\end{enumerate}


\section{Protesty, odwołania}
\begin{enumerate}
    \item Protest to zgłoszenie przez załogę (jacht) lub Sędziego zarzutu naruszenia prawa drogi lub niniejszego Regulaminu przez inną załogę (jacht) podczas Regat.
    \item Protestujący jacht musi poinformować jacht protestowany jak najszybciej po incydencie -- wyraźny okrzyk \textit{,,Protest''}.
    Jeżeli nie jest w zasięgu słyszalności, jacht prostestujący musi poinformować drugi przy pierwszej okazji.
    \item Sędzia prostestujący jacht musi poinformować jacht protestowany przy pierwszej okazji, najpóźniej do 15 minut po zakończeniu biegu.
    \item Protest musi określać (najlepiej w formie pisemnej):
    \begin{enumerate}
        \item Dane protestującego i protestowanego
        \item Incydent (opis sytuacji, czas, miejsce)
        \item Możliwy naruszony przepis, który w opinii protestującego został naruszony
    \end{enumerate}
    \item Protest należy złożyć Sędziemu w ciągu 30 minut od zakończenia biegu, w którym nastąpił incydent.
    \item Sędzia określa czas i miejsce rozpatrzenia protestu, które powinno nastąpić jak najszybciej, w szczególności przed ogłoszeniem wyników danej \textit{Części}.
    \item Na rozpatrzeniu protestu obowiązkowa jest obecność przedstawicieli protestującego i protestowanego jachtu.
    Prostestujący jest odpowiedzialny za obecność przedstawiciela jachtu protestowanego, bez tego protest może zostać odrzucony.
\end{enumerate}


\section{Kary, wykluczenia}
\begin{enumerate}
    \item Za przekroczenie dopuszczalnego czasu biegu następuje dyskwalifikacja pomiaru (w szczególnych przypadkach dopuszczalne jest 15 sekund nadużycia czasu).
    \item Kara obrotów -- odbywa się po jak najszybszym uwolnieniu się od pozostałych jachtów, przed przepłynięciem mety. Polega na wykonaniu odpowiedniej liczby obrotów (360°), w jednym kierunku. Jeden obrót zawiera jeden zwrot na wiatr (przez sztag) i jeden z wiatrem (przez rufę).
    \item Za wymuszenie prawa drogi:
    \begin{enumerate}
        \item Jacht, który wymusił prawo drogi na innym jachcie, winien wykonać karę \textbf{dwóch obrotów}
        \item W przypadku niewykonania kary obrotów podczas biegu, następuje pogorszenie otrzymanego wyniku o $15\%$ -- odpowiednio zwiększenie dla czasu, zmniejszenie dla dystansu.
    \end{enumerate}
    \item Za dotknięcie znaku (boi):
    \begin{enumerate}
        \item Jacht, który dotknął znaku, winien wykonać karę \textbf{jednego obrotu}
        \item W przypadku niewykonania kary obrotów podczas biegu, następuje pogorszenie otrzymanego wyniku o $8\%$ -- odpowiednio zwiększenie dla czasu, zmniejszenie dla dystansu.
    \end{enumerate}
    \item W przypadku znaczącego pogorszenia wyniku (lub miejsca) jachtu nie z jego winy (np. udzielanie pomocy innemu jachtowi), przysługuje prośba o zadośćuczynienie do Sędziego, który może przyznać odpowiednią rekompensatę po rozpatrzeniu wniosku.
    Prośby można składać na Odprawie po biegu, przed ogłoszeniem wyników danej \textit{Części} Regat.
    \item Nieprzybycie do pomostu w określonym przez Sędziego czasie skutkuje pogorszeniem otrzymanego wyniku proporcjonalnie do długości spóźnienia. Za każdą minutę spóźnienia wynik zmniejsza się o $1\%$.
    Spóźnienie powyżej 15 minut, uniemożliwiające terminowe przeprowadzenie kolejnych biegów, skutkuje dyskwalifikacją z Regat.
\end{enumerate}

\section{Postanowienia końcowe}
\begin{enumerate}
    \item Organizator zastrzega sobie prawo do zmiany Regulaminu, jedynie przed rozpoczęciem Regat. W szczególności może to dotyczyć zasad punktacji, ograniczenia formatu regat do samej Formuły Smartfonowej, np. przy mniejszej liczbie załóg.
    \item Wszelkie zmiany w Regulaminie muszą być ogłoszone najpóźniej przed zapisami, odbywającymi się przed rozpoczęciem Regat.
\end{enumerate}


\newpage
\appendix

% Wyśrodkowanie tytułów sekcji, bez paragrafu
\titleformat{\section}
  {\normalfont\large\bfseries\centering}{\thesection.}{1em}{}

\section{Regatowe Prawo Drogi (RRS) -- Skrót}
\label{sec:RegatowePrawoDrogi}
% TODO
\begin{enumerate}
    \item Skrót przygotowany na podstawie Przepisów Regatowych Żeglarstwa PZŻ i Recing Rules of Sailing 2021-2024.
    \item Prawa podstawowe:
    \begin{enumerate}
        \item 
    \end{enumerate}
    \item Ograniczenia ogólne:
    \begin{enumerate}
        \item aa
    \end{enumerate}
    \item W pobliżu znaków i przeszkód:
    \begin{enumerate}
        \item 
    \end{enumerate}
\end{enumerate}

\end{document}

